\documentclass[11pt]{article}
\usepackage[margin=1in]{geometry}
\usepackage{amsmath,amssymb}
\usepackage{graphicx}
\usepackage{booktabs}
\usepackage[hidelinks]{hyperref}
\graphicspath{{figures/}}

\title{Research Project: Working Report}
\author{Research team}
\date{\today}

\begin{document}
\maketitle

\begin{abstract}
This document is a report scaffold. No scientific findings have yet been
entered. Reviewed findings will be added with their scope, uncertainty,
citations and computational provenance.
\end{abstract}

\section{Research objective}
Define the question, motivation and scope after project initialization.

\section{System and theoretical framework}
Define observables, notation, assumptions, governing equations,
approximations and their domains of applicability.

\section{Methods and validation}
Describe analytical, computational or experimental methods, inputs,
convergence checks, controls and uncertainty estimates.

\section{Results}
Insert reviewed findings only. Each computational claim should have a
nearby source comment identifying its result ID and evidence manifest.
% Example format only: R-ID; results/<run-id>/manifest.json; observable/data locator.

\section{Discussion and limitations}
Separate observations from interpretations. Discuss alternatives,
applicable regimes and unresolved discrepancies explicitly.

\section{Conclusions and next questions}
Summarize approved conclusions with their limits and identify open questions.

% Add verified citation entries and \cite{key} commands when sources are used.
% No placeholder scientific citation is inserted in this scaffold.
\bibliographystyle{plain}
\bibliography{references}

\end{document}